% chapters/08/rangkuman.tex
\bagian{Rangkuman}

\begin{rangkumanbox}
\begin{itemize}
  \item Pada bab ini benda menjadi agen: kecepatannya tidak lagi berasal dari
        tombol atau dari gravitasi, melainkan dari keputusannya sendiri.
  \item Susunan keputusan sebuah agen terdiri atas tiga langkah, yaitu
        persepsi, keputusan, dan tindakan. Susunan itu tetap sama dari agen
        yang paling sederhana sampai yang paling rumit.
  \item Persepsi menghasilkan dua hal dari satu perhitungan: vektor menuju
        pemain dan jaraknya. Jarak dipakai untuk memutuskan, arah dipakai untuk
        bergerak.
  \item Jari-jari persepsi membatasi jangkauan pengamatan NPC. Nilainya
        merupakan pilihan perancang, bukan hasil perhitungan.
  \item Kecepatan dihitung dengan menormalkan vektor arah lalu mengalikannya
        dengan kecepatan yang ditetapkan, sehingga besarnya tetap berapa pun
        jaraknya.
  \item Memakai vektor jarak secara langsung sebagai kecepatan merupakan cara
        yang salah, sebab kecepatannya menjadi sama dengan jaraknya. Pengukuran
        pada program memperlihatkan kecepatan tetap 150,0 piksel per detik pada
        cara yang benar, dan 75,2 sampai 184,0 piksel per detik pada cara yang
        salah.
  \item Jarak henti diperlukan agar NPC tidak bergetar di atas pemain. Jarak
        henti harus lebih besar daripada perpindahan satu bingkai, yaitu
        kecepatan dibagi laju bingkai.
  \item \istilah{Hysteresis} pada batas persepsi---memakai dua nilai batas yang
        berbeda untuk mulai dan berhenti mengejar---mencegah pergantian
        keputusan yang berlebihan ketika pemain berada tepat pada batas.
  \item AI reaktif hanya memakai keadaan saat itu, sehingga murah dan mudah
        diperiksa, tetapi tidak mengetahui adanya rintangan. AI berbasis
        perencanaan memperhitungkan lebih banyak keterangan, dengan biaya
        perhitungan yang lebih besar.
  \item NPC pada bab ini tidak melakukan kesalahan ketika menekan dinding; yang
        kurang adalah keterangan yang tersedia baginya. Menambah keterangan
        itulah yang menjadi pekerjaan Bab~IX dan Bab~X.
\end{itemize}
\end{rangkumanbox}
